% A alocação eficiente de recursos para a execução de processos é fundamental em diferentes indústrias de transporte, distribuição e produção. Podemos modelar esses processos de alocação como problemas de otimização. Um caso específico desses problemas é o \textit{Just-In-Time Job Shop Scheduling} (JITJSS), no qual as tarefas a serem executadas possuem prazos de entrega, e finalizar essas tarefas fora do seu prazo gera penalidades de adiantamento e atraso. Com isso, nosso objetivo é minimizar a soma dessas penalidades.
% Sabemos que o JITJSS é um problema fortemente NP-difícil, isto é, uma solução exata é computacionalmente inviável na prática. Por isso, uma alternativa para resolver problemas desse tipo é utilizar heurísticas, ou seja, estratégias que obtêm soluções em um tempo viável, mas sem garantia de qualidade. Porém, a construção de heurísticas gera, para o projetista, a dificuldade de testar manualmente diferentes combinações de estratégias para encontrar uma combinação que componha uma heurística eficiente. Por isso, o processo se torna tedioso, custoso, dificilmente reproduzível e inconclusivo.
% Para mitigar esses problemas, a comunidade científica desenvolveu diversas ferramentas e técnicas para a Configuração Automática de Algoritmos (do inglês, \textit{Automatic Algorithm Configuration} -- AAC) e Projeto Automático de Algoritmos (do inglês, \textit{Automatic Algorithm Design} -- AAD). Essas ferramentas e técnicas sistematizam e automatizam a parte mecânica do desenvolvimento de heurísticas, permitindo ao projetista criar algoritmos mais eficientes. De fato, a eficácia das técnicas de AAC e AAD na geração de algoritmos de nível estado da arte já está comprovada pelos resultados de diversos trabalhos na literatura. Porém, a aplicação dessas técnicas ao JITJSS ainda não foi plenamente explorada. A partir disso, o objetivo deste trabalho é aplicar a AAC e o AAD ao JITJSS, propondo um novo espaço de projeto algorítmico e gerando, de forma automática, novos algoritmos para o problema. Esperamos obter soluções que superem o atual estado da arte e validem a efetividade dos componentes construídos. 

% \keywords{Projeto Automático de Algoritmos, Configuração Automática de Algoritmos, Problemas de Agendamento, Heurísticas}

A alocação eficiente de recursos para a execução de processos é fundamental em diferentes indústrias de transporte, distribuição e produção. Podemos modelar esses processos de alocação como problemas de otimização. Um caso específico desses problemas é o \textit{Just-In-Time Job Shop Scheduling} (JITJSS), no qual as tarefas a serem executadas possuem prazos de entrega, e finalizar essas tarefas fora do seu prazo gera penalidades de adiantamento e atraso. Com isso, nosso objetivo é minimizar a soma dessas penalidades. Sabendo que o JITJSS é um problema fortemente NP-difícil, uma solução exata é computacionalmente inviável na prática. Por isso, uma alternativa para resolver problemas desse tipo é utilizar heurísticas, ou seja, estratégias que obtêm soluções em um tempo viável, mas sem garantia de qualidade. Porém, a construção de heurísticas gera, para o projetista, a dificuldade de testar manualmente diferentes combinações de estratégias para encontrar uma combinação que componha uma heurística eficiente. Por isso, o processo se torna tedioso, custoso, dificilmente reproduzível e inconclusivo. Para mitigar esses problemas, a comunidade científica desenvolveu diversas ferramentas e técnicas para a Configuração Automática de Algoritmos (do inglês, \textit{Automatic Algorithm Configuration} -- AAC) e Projeto Automático de Algoritmos (do inglês, \textit{Automatic Algorithm Design} -- AAD). Essas ferramentas e técnicas sistematizam e automatizam a parte mecânica do desenvolvimento de heurísticas, permitindo ao projetista criar algoritmos mais eficientes. De fato, a eficácia das técnicas de AAC e AAD na geração de algoritmos de nível estado da arte já está comprovada pelos resultados de diversos trabalhos na literatura. Porém, a aplicação dessas técnicas ao JITJSS ainda não foi plenamente explorada. A partir disso, o objetivo deste trabalho é aplicar a AAC e o AAD ao JITJSS, propondo um novo espaço de projeto algorítmico e gerando, de forma automática, novos algoritmos para o problema. Para tanto, estruturou-se um espaço de busca parametrizado via gramática livre de contexto. O processo de configuração convergiu para uma arquitetura baseada em Busca Tabu Iterada, validando a relevância da vizinhança original proposta, \texttt{troca\_crit}, que esteve presente em 80\% das soluções elites. Nos testes experimentais realizados sobre 72 instâncias de teste da literatura, o algoritmo gerado igualou ou superou o melhor resultado conhecido em aproximadamente 72,2\% dos cenários. Ademais, o método demonstrou alta eficiência temporal, localizando as soluções finais em uma média de 21,1 segundos. Os resultados obtidos validam a efetividade dos componentes construídos e chancelam o projeto automático como uma alternativa robusta à calibração manual de heurísticas para problemas complexos de agendamento.

\keywords{Projeto Automático de Algoritmos, Configuração Automática de Algoritmos, Problemas de Agendamento, Heurísticas}